
\subsection{Introduction}
%
The objective of this section is to develop a framework for formulating constitutive equations for the models introduced in \cref{sec:genwarp}, which relates the generalized strains \(\bm{e}\) and stresses \(\bm{s}\):
\RepeatEquation{eq:cosserat-strain}

\subsubsection{Literature}

\subsubsection{Objective}

\subsubsection{Outline}

\Cref{sec:enhan-prelim} introduces the reference continuum model and its connection to the 1D fields of the beam theory \eqref{eq:finite-alignment-field}. 
The constrained stress condition \eqref{eq:saint-venant-stress} is introduced and justified. 

\Cref{sec:section-kinematics} derives the \emph{compatible} strain field in the setting of \cref{sec:enhan-prelim}, where a natural distinction arises between geometric nonlinearities at the \emph{section} and \emph{element} levels. 

\Cref{sec:enhan-section} introduces the enhanced strain principle. 


\subsection{Preliminaries}\label{sec:enhan-prelim}

In the continuum setting the model is a body %with reference configuration 
\(
\Body=(0,L)\times\Section
\),
where \(\Section\subset\mathbb{R}^2\) is a bounded cross-section
spanned by basis vectors \(\{\FrameBasis_y,\FrameBasis_z\}\) in \(\mathbb{R}^3\). 
Together with the axial unit vector \(\FrameBasis\), the vectors \(\{\FrameBasis,\,\FrameBasis_y,\,\FrameBasis_z\}\) form an orthonormal basis and the 3D identity is:
\[
\mathbf{1}=\FrameBasis\otimes\FrameBasis+\oneS
\qquad\text{with}\qquad
\oneS=\FrameBasis_{\beta}\otimes\FrameBasis_{\beta}
\]
A point in \(\Body\) has coordinates \(\reflenx\,\FrameBasis + \bm{r}\) in the reference configuration, where \(\bm{r} \in\Section_x\) is the coordinate of a point in the section at \(\reflenx\) and \(\bm{r}\cdot\FrameBasis = 0\).
The boundary of \(\Body\) consists of the the lateral boundary \(\BodyLateralSurface=\SectionBoundary\times(0,L)\) and the surfaces \(\Section_0,\,\Section_L\) of the sections at the ends \(\reflenx \in \{0,L\}\).  

The deformation is represented by the deflection \(\bm{x}(x)\in\mathbb{R}^3\), rotation \(\FrameRotation(x)\in\mathrm{SO}(3)\) and warping fields \(\bm{\alpha}(x)\in\mathbb{R}^{n_w}\):
\begin{SaveEquation}{eq:finite-alignment-field}
\hat{\bm{x}}(\reflenx, \PlanePosition) 
= \bm{x}(\reflenx) + \FrameRotation \bigl(\PlanePosition 
+ \WarpTotalShapes(\PlanePosition)\bm{\alpha}(\reflenx)\bigr)
% \label{eq:finite-alignment-field}
\end{SaveEquation}
%
where \(\bm{\Phi}:\Section\rightarrow\mathbb{R}^3\) represents cross sectional warping shapes.  
The representation \eqref{eq:finite-alignment-field} is exact in the limit \(n_w\rightarrow\infty\). 

\begin{remark}
The linearization of \eqref{eq:finite-alignment-field} is:
%
\begin{equation}
\hat{\bm x}(\reflenx, \bm{r})=(\reflenx\FrameBasis+\bm{r})+\bm\theta(\reflenx)\times\bm{r}+\bm\Phi(\bm{r})\,\bm\alpha(\reflenx)
\label{eq:linear-alignment-field}
\end{equation}
and the linearized displacement \(\hat{\bm{u}} = \hat{\bm{x}} - (x\FrameBasis + \bm{r})\) is
\begin{SaveEquation}{eq:trace-lift-linear}
\hat{\bm u}(\reflenx, \bm{r})=\bm{u}(x)
+\bm\theta(x)\times\bm{r}+\bm\Phi(\bm{r})\,\bm{\alpha}(x)
% \label{eq:trace-lift-linear}
\end{SaveEquation}
\end{remark}


\begin{remark}
The following developments are made without appealing to a coordinate system (we only assume right-handedness), and thus at no commitments are made to a particular coordinate convention. However, such conventions can aid the interpretation of equations and are required for implementation. 
Two common conventions for assigning coordinates \(x,y,z\) are:
\begin{itemize}
    \item Take the first (\(x\)) coordinate to coincide with the reference longitudinal axis \(\FrameBasis\). This furnishes components on the basis \(\{\bm{\partial}_x,\bm{\partial}_y,\bm{\partial}_z\}\)
    \[
    \FrameBasis \equiv \bm{\partial}_x = \begin{pmatrix}
        1 \\ 0 \\ 0
    \end{pmatrix},
    \quad
    \PlanePosition = \begin{pmatrix}
    0 \\ y \\ z
    \end{pmatrix}
    \qquad\text{and}\qquad 
    \FrameBasis^{\times} = \begin{bmatrix}
        0 & 0 & 0 \\
        0 & 0 & -1 \\
        0 & 1 & 0
    \end{bmatrix}
    \]
    \item Take the last (\(z\)) coordinate to coincide with \(\FrameBasis\). This furnishes components on the basis \(\{\bm{\partial}_x,\bm{\partial}_y,\bm{\partial}_z\}\)
    
    \[
    \FrameBasis \equiv \bm{\partial}_z = \begin{pmatrix}
        0 \\ 0 \\ 1
    \end{pmatrix},
    \quad
    \PlanePosition = \begin{pmatrix}
    x \\ y \\ 0
    \end{pmatrix}
    \qquad\text{and}\qquad 
    \FrameBasis^{\times} = \begin{bmatrix}
        0 & -1 & 0 \\
        1 &  0 & 0 \\
        0 & 0 & 0
    \end{bmatrix}
    \]
\end{itemize}
\end{remark}
\begin{remark}
The 3D out-of-plane reference director \(\FrameBasis\) is used to form cross products with vectors in the plane \(\Section\). 
For example, given a function \(u: S \rightarrow \mathbb{R}\) with \(u' = 0\), the notation \(\FrameBasis \times \nabla u\) may be interpreted in coordinates \(\bm{\partial}_x \coloneq \FrameBasis\) either as 
\[
\begin{bmatrix}
0 & 0 & 0 \\
0 & 0 & -1 \\
0 & 1 & 0
\end{bmatrix}
\begin{pmatrix}
0 \\ u_{,y} \\ u_{,z}
\end{pmatrix}
\qquad\text{or}\qquad 
\begin{bmatrix}
0 & -1 \\
1 & 0
\end{bmatrix}
\begin{pmatrix}
u_{,y} \\ u_{,z}
\end{pmatrix}
\]
\end{remark}

\subsubsection{Assumptions}


% \subsubsection{Assumptions} %Stress Constraints}
% Recall that the Saint-Venant field generates vanishing stress components \(\FrameBasis_{\alpha}\cdot\PlaneStressTensor\FrameBasis_{\beta} = 0\). 
The stress tensor is assumed to take the form:
%
\begin{SaveEquation}{eq:saint-venant-stress}
\StressTensor = \AxialStress\, \FrameBasis \otimes \FrameBasis+\ShearStress \otimes \FrameBasis+\FrameBasis \otimes \ShearStress
\qquad\text{with}\qquad
\begin{aligned}
\AxialStress    &\coloneq \FrameBasis\cdot \StressTensor\, \FrameBasis \\
\ShearStress &\coloneq \FrameBasis_{\beta}\otimes \FrameBasis_{\beta} \StressTensor\,\FrameBasis
\end{aligned}
% \label{eq:frame-stress-tensor}
% \label{eq:saint-venant-stress}
\end{SaveEquation}
where \(\AxialStress\) is the scalar axial stress and \(\ShearStress \in \mathbb{R}^2\) the vector of shear stresses. 
This furnishes the representation for the stress vector \(\StressVector \in \mathbb{R}^3\):
\begin{equation}
\StressVector := \StressTensor\,\FrameBasis = \AxialStress \FrameBasis + \ShearStress
\label{eq:def-stress-vector}
\end{equation}
The infinitesimal strain tensor $\StrainTensor$ is decomposed as 
\begin{SaveEquation}{eq:strain-tensor-decomp}
\StrainTensor = \FrameBasis\stackrel{s}{\otimes}(\PointAxialStrain\,\FrameBasis 
+ \PointShearStrain) + \PlaneStrainTensor_{\alpha\beta}\FrameBasis_{\alpha}\otimes\FrameBasis_{\beta} 
% \label{eq:strain-tensor-decomp}
\end{SaveEquation}
with $\PointShearStrain$ a vector with zero projection on $\FrameBasis$. 
For any superposition satisfying \(\StressTensor\,\bm n=\mathbf{0}\) on \(\BodyLateralSurface\) and \eqref{eq:saint-venant-stress},
the field equations reduce on each section to
\begin{subequations}
\label{eq:sv-equilibrium}
\begin{align}
\bm{\tau}'&=\mathbf{0} && \text{in }\Section,\\
\DivS\bm{\tau}+\AxialStress'&=0 && \text{in }\Section, \label{eq:sv-axial-equilibrium}\\
\bm{\tau}\cdot\bm\nu&=0 && \text{on }\SectionBoundary,
\end{align}
\end{subequations}
where the prime denotes derivative with respect to \(\reflenx\).

\iffalse
% \paragraph{Isotopy}

Similarly decompose the stress as 
\[
\StressTensor = \AxialStress \FrameBasis\otimes\FrameBasis
+ \ShearStress\otimes\FrameBasis
+ \FrameBasis\otimes\ShearStress
+ \PlaneStressTensor
\]
with $\ShearStress$ a vector satisfying $\FrameBasis\cdot\ShearStress = 0$ and $\PlaneStressTensor = \PlaneStressTensor_{\alpha\beta}\FrameBasis_{\alpha}\otimes\FrameBasis_{\beta}$. 
Under isotropy, the shear stress \(\ShearStress\) is coupled only to the shear strain \(\bm{\varepsilon}_{\gamma}\):
\[
\bm{\tau} = G\bm{\varepsilon}_{\gamma}
\]
However, the axial \(\AxialStress\) and section \(\hat{\StressTensor}\) stresses are coupled:
\[
\begin{aligned}
\hat{\StressTensor}
% &=2 G \hat{\bm{E}}+\lambda(\varepsilon+\operatorname{tr} \hat{\bm{E}})\left(\mathrm{i}_y \otimes \FrameBasis_y+\FrameBasis_z \otimes \mathrm{i}_z\right)\\
% &=\frac{E}{1+\nu} \hat{\bm{E}}+\frac{E \nu}{(1+\nu)(1-2 \nu)}(\varepsilon+\operatorname{tr} \hat{\bm{E}}) \; \mathrm{i}_{\alpha} \otimes \mathrm{i}_{\alpha}\\
&= 2G \hat{\bm{E}}+\lambda (\PointAxialStrain + \operatorname{tr} \hat{\bm{E}}) \; \FrameBasis_{\alpha} \otimes \FrameBasis_{\alpha} \\
% \end{aligned}
% \]
% \[
% \begin{aligned}
\AxialStress
&=(2 G+\lambda) \PointAxialStrain+\lambda \operatorname{tr} \PlaneStrainTensor \\
% &=2 G \frac{1-\nu}{1-2 \nu} \varepsilon+\frac{2 G \nu}{1-2 \nu} \operatorname{tr} \hat{\bm{E}} \\
% &=\frac{E(1-\nu)}{(1+\nu)(1-2 \nu)} \varepsilon+\frac{E \nu}{(1+\nu)(1-2 \nu)} \operatorname{tr} \hat{\bm{E}} \\
% &= M\varepsilon + \lambda \operatorname{tr}\hat{\StrainTensor}
\end{aligned}
\]
% with \(M=\frac{E(1-\nu)}{(1+\nu)(1-2 \nu)}\) the \emph{P-wave modulus}. 
Thus, if the orthogonal strain \(\PlaneStrainTensor\) satisfies:
\begin{equation*}
\operatorname{tr}\PlaneStrainTensor = -2\nu\,\PointAxialStrain
\quad\text{then}\quad 
\AxialStress = E\,\PointAxialStrain
\end{equation*}


This furnishes the inner product equalities
\begin{equation*}
\left<\TraceTotalEnergyStress,\,\delta\FrameStrain\right>
\equiv \left<\StressVector,\, \delta \hat{\StrainVector}\right>
\equiv \left<\StressTensor ,\,\delta\StrainTensor\right>
\end{equation*}

For Saint–Venant-type stresses with the form \eqref{eq:saint-venant-stress} one has
\begin{equation}
\left<\StressTensor,\,\delta\hat{\bm{E}}\right>
=\left<\StressVector,\,\delta\hat{\bm{\varepsilon}}\right>,
\quad\text{since}\quad
\left<\StressTensor,\,\stackrel{\rm s}{\nabla}\PlaneWarp\,\right>=0,
\label{eq:sv-orthogonality}
\end{equation}

\begin{ambox}[Material State Determination]{box:mati}
\textbf{Objective.} 
Given a strain vector \(\StrainVector \in \mathbb{R}^3\) and tensorial stress-strain material model \(\TensorMaterial \colon \mathrm{Sym}(\mathbb{R}^3, \mathbb{R}^3) \rightarrow \mathrm{Sym}(\mathbb{R}^3, \mathbb{R}^3)\), return the stress vector \(\StressVector \in \mathbb{R}^3\) representing the constrained stress tensor.
% \[
% \StressTensor = \StressVector \symbun \FrameBasis
% \]
\textbf{Procedure.} 
Perform Newton-Raphson iterations solving for the strain tensor \({\PlaneStrainTensor}\) that satisfies the constraints
\[
\PlaneBasis_{\alpha}\cdot  \TensorMaterial\bigl(\StrainVector \symbun \FrameBasis + \PlaneStrainTensor \bigr)\, \PlaneBasis_{\beta} = 0
\]
Details of the procedure are developed by \cite{klinkel2002using}.
\end{ambox}

\fi 

\iffalse
\begin{Quote}
Recall
\(
\FrameBasis\cdot\PlaneWarp=0.
\)
hence
\[
\FrameBasis\cdot\partial_{\beta}\PlaneWarp=\partial_{\beta}(\FrameBasis\cdot\PlaneWarp)=0
\quad\Rightarrow\quad
\big(\nabla\PlaneWarp\big)\FrameBasis=\mathbf{0},
\qquad
\FrameBasis\cdot\big(\nabla\PlaneWarp\big)=\mathbf{0}^{\top}.
\]
where \(\nabla(\cdot)\) acts on section coordinates \(\bm{r}=(y,z)\) only.
Therefore \(\stackrel{\rm s}{\nabla}\PlaneWarp\,\) has \emph{only} in–plane components \(\FrameBasis_\alpha\otimes\FrameBasis_\beta\), with
\[
\big[\sym (\nabla\PlaneWarp)\big]_{\alpha i}
=\tfrac12\big((\pi_{\alpha,i}+\pi_{i,\alpha}\big)
=\tfrac12\big(0+0\big)=0,
\qquad
\big[\stackrel{\rm s}{\nabla}\PlaneWarp\,\big]_{ii}=0.
\]

Now assume the Saint–Venant stress structure \eqref{eq:saint-venant-stress} so that \(\StressTensor_{\alpha\beta}=0\). 
Then, pointwise,
\[
\StressTensor:\stackrel{\rm s}{\nabla}\PlaneWarp\,
=\underbrace{\sigma\,\FrameBasis\otimes\FrameBasis:\stackrel{\rm s}{\nabla}\PlaneWarp\,}_{\displaystyle \sigma\,\FrameBasis\cdot\stackrel{\rm s}{\nabla}\PlaneWarp\,\,\FrameBasis=0}
+\underbrace{\bm{\tau}\otimes\FrameBasis \cdot \stackrel{\rm s}{\nabla}\PlaneWarp\,}_{\displaystyle \bm{\tau}\cdot\stackrel{\rm s}{\nabla}\PlaneWarp\,\,\FrameBasis=0}
+\underbrace{\FrameBasis\otimes\bm{\tau} \cdot \stackrel{\rm s}{\nabla}\PlaneWarp\,}_{\displaystyle \FrameBasis\cdot\stackrel{\rm s}{\nabla}\PlaneWarp\,\,\bm{\tau}=0}
=0,
\]
because \(\stackrel{\rm s}{\nabla}\PlaneWarp\,\,\FrameBasis=\mathbf{0}\) and \(\FrameBasis \cdot\stackrel{\rm s}{\nabla}\PlaneWarp\,=\mathbf{0}\). Consequently,
\[
\left<\StressTensor,\,\stackrel{\rm s}{\nabla}\PlaneWarp\right>=\SectionIntegral{\StressTensor \cdot \stackrel{\rm s}{\nabla}\PlaneWarp\,}=0,
\]
and the term is orthogonal to the Saint–Venant stress subspace \(\{\sigma\,\FrameBasis\otimes\FrameBasis+\bm{\tau}\otimes\FrameBasis+\FrameBasis\otimes\bm{\tau}\}\).

The rod shear components are the mixed \(\FrameBasis_\alpha\otimes\FrameBasis\) entries of \(\hat{\bm{E}}\), i.e.\ those generated by
\[
\hat{\bm{\varepsilon}}
=\bm{\gamma}
+\bm{\kappa}\times\bm{r}
+ \bm{\kappa}\times\bm{\Phi}\bm{\alpha}
+ \bm{\Phi}\,\bm{\alpha}'
+ \nabla\omega,
\qquad
\omega\coloneq \FrameBasis\cdot(\bm{\Phi}\bm{\alpha}),
\]
since \(\hat{\bm{E}}=(\hat{\bm{\varepsilon}}\stackrel{s}{\otimes}\FrameBasis)+\underbrace{\stackrel{\rm s}{\nabla}\PlaneWarp\,}_{\alpha\beta\text{ only}}+\mathrm{h.o.t.}\).
Only \(\hat{\bm{\varepsilon}}\) carries \(\FrameBasis_\alpha\otimes\FrameBasis\) (rod shear) and axial \(\FrameBasis\otimes\FrameBasis\)  components and hence couples to \(\StressVector=\StressTensor\,\FrameBasis\).

Thus \(\stackrel{\rm s}{\nabla}\PlaneWarp\,\) may create in–plane (membrane) strain/shear \(\alpha\beta\), but not the mixed \(\alpha i\) rod shear, and it therefore produces no virtual work against \(\StressTensor\) of Saint–Venant form.
\end{Quote}
\fi


\subsection{Compatible Strain}\label{sec:section-kinematics}

\paragraph{Finite Strain.}
The tangent map \(\bm{F}\coloneq T\,\hat{\chi}\) of the finite deformation \eqref{eq:finite-alignment-field} is:
\begin{equation}
\bm{F}
=\FrameRotation\!\left(\mathbf{1}+\bm{\mu}\otimes\FrameBasis
~+~\FrameBasis\otimes\nabla\omega
~+~\nabla\PlaneWarp\right),
\qquad\text{with}\qquad
\bm{\mu}(\bm{r})
=\bm{\gamma}+\bm{\kappa}\times\bigl(\bm{r}+\bm{\Phi}\bm{\alpha}\bigr)
+\bm{\Phi}\,\bm{\alpha}'
\label{eq:F-Phi}
\end{equation}
where $\ShearStrain $ and $\CurveStrain $ are the beam strain variables defined by \eqref{eq:def-gam-kap}, and the warping is decomposed into axial and transverse parts
\begin{equation*}
\AxialWarp\coloneq \FrameBasis\cdot\bm{\Phi}\bm{\alpha}\in\mathbb{R}, 
\qquad\text{and}\qquad
% \PlaneWarp\coloneq \mathbf{1}_{\Section}\,\bm{\omega}
\PlaneWarp\coloneq \mathbf{1}_{\Section}\,\bm{\Phi}\bm{\alpha}
,
\end{equation*}
Introduce the decomposition:
\begin{equation}
\WarpTotalShapes = \FrameBasis\otimes\WarpAxialShapes + \WarpPlaneShapes,
\qquad
\WarpAxialShapes\coloneq \FrameBasis\cdot\bm{\Phi}\in\mathbb{R}^{1\times n_w}.
\label{eq:decom-warp-shapes}
\end{equation}

The Green-Lagrange strain tensor \(\hat{\StrainTensor}\coloneq\tfrac12(\bm{F}^{\top}\bm{F}-\mathbf{1})\) with \eqref{eq:F-Phi} is:
\begin{equation}
\begin{aligned}
\hat{\StrainTensor}
&=~\bm{\mu}\stackrel{\mathrm{s}}{\otimes}\FrameBasis
~+~\nabla\AxialWarp\stackrel{\mathrm{s}}{\otimes}\FrameBasis
~+\stackrel{\rm s}{\nabla}\PlaneWarp
\\
&\quad
+~(\bm{\mu}\cdot\FrameBasis)\,\nabla\AxialWarp\stackrel{\mathrm{s}}{\otimes}\FrameBasis
~+~\tfrac12\Big(
\bm{\mu}\cdot\bm{\mu}\,\FrameBasis\otimes\FrameBasis
+ \nabla\AxialWarp\otimes\nabla\AxialWarp
+ (\nabla\PlaneWarp)^{\top}\nabla\PlaneWarp
\Big).
\end{aligned}
\label{eq:exact-gl-strain-Phi}
\end{equation}
where the symmetric bun product \(\bm{a}\stackrel{\mathrm{s}}{\otimes}\bm{b}\) is defined for vectors \(\bm{a}\), \(\bm{b}\), and \(\bm{c}\) by
\begin{equation*}
    (\bm{a}\stackrel{\mathrm{s}}{\otimes}\bm{b} )\bm{c} \coloneq \tfrac{1}{2} \bigl(
      \bm{a}\,(\bm{b}\cdot\bm{c}) + \bm{b}\,(\bm{a}\cdot\bm{c})
    \bigr)
\end{equation*}


\paragraph{Moderate (Wagner) Strain.}
A truncation of \eqref{eq:exact-gl-strain-Phi} which retains important nonlinearities is:
\begin{SaveEquation}{eq:wagner-gl-strain-pi}
\hat{\StrainTensor}
~=~\hat{\StrainVector}\stackrel{\mathrm{s}}{\otimes}\FrameBasis
~+~\stackrel{\rm s}{\nabla}\PlaneWarp
~+~\mathrm{h.o.t.},
\qquad
%
\hat{\StrainVector}\coloneq
\bm{\gamma}
+\bm{\kappa}\times\bm{r}
+\bm{\Phi}\,\bm{\alpha}'
+\nabla\AxialWarp
+\tfrac{1}{2}\Twist^2\,
\bm{r}\cdot\bm{r}\,\FrameBasis
% \label{eq:wagner-gl-strain-pi}
\end{SaveEquation}
Linearizing \cref{eq:wagner-gl-strain-pi} furnishes
\[
\delta \hat{\bm{E}} = \delta \ShearStrain - \bm{r}\times\delta \CurveStrain + (\FrameBasis\otimes \bm{\varphi})\delta \bm{\alpha} + (\nabla \bm{\varphi})^{\top} \delta \bm{\alpha}' + \Twist r^2\, \FrameBasis\otimes\FrameBasis\delta \Twist 
\]
where \(r^2 \coloneq \bm{r}\cdot\bm{r}\).
\iftrue
From \eqref{eq:wagner-gl-strain-pi} the  variation
\(
\delta\StrainVector=\TotalStrainMatrix\,\delta\bm{e}
\)
furnishes the section kinematic matrix:
\begin{equation}
\TotalStrainMatrix
=\left[
\;\mathbf{1}
~~~-\,\bm{r}^{\times}+\Twist r^2 \FrameBasis\otimes\FrameBasis
~~~~\FrameBasis\otimes\bm{\varphi}
~~~~\FrameBasis_{\beta}\otimes\nabla_{\beta}\WarpAxialShapes
% ~~~\Big(\,\bm{\kappa}^{\times}\WarpTotalShapes~+~\nabla\WarpAxialShapes\,\Big)
\right],
\label{eq:def-section-kinematic-matrix-wagner}
\end{equation}
where 
\(
(\nabla\WarpAxialShapes)\,\delta\bm{\alpha}
=\nabla\big((\FrameBasis^{\top} \bm{\Phi})\,\delta\bm{\alpha}\big)
=\nabla\omega\big|_{\delta\bm{\alpha}}
\),
and \(\gradS\) acts with respect to the section coordinates \(\bm{r}\).
Thus, for any virtual increment,
\begin{equation}
\delta\hat{\bm{\varepsilon}}
=\delta\bm{\gamma}
-\big(\bm{r}
+\bm{\Phi}\bm{\alpha}\big)\times\delta\bm{\kappa}
+\bm{\Phi}\,\delta\bm{\alpha}'
+\bm{\kappa}\times\bm{\Phi}\,\delta\bm{\alpha}
+\nabla\big(\FrameBasis\cdot\bm{\Phi}\,\delta\bm{\alpha}\big),
\label{eq:delta-strain-vector-wagner}
\end{equation}
and the remainder
\(
\delta\hat{\bm{E}}_{\perp}=\sym\big(\nabla(\mathbf{1}_{\Section}\bm{\Phi}\,\delta\bm{\alpha})\big)
\)
is orthogonal to \(\StressTensor\) under the constrained stress condition \eqref{eq:saint-venant-stress}.
\fi

\paragraph{Linear Strain.}
The first order truncation of \eqref{eq:exact-gl-strain-Phi} furnishes:
\begin{SaveEquation}{eq:linear-gl-strain-pi}
% \StrainTensor
% ~=~\StrainVector\stackrel{\mathrm{s}}{\otimes}\FrameBasis
% ~+~\stackrel{\rm s}{\nabla}\PlaneWarp
% \qquad
%
\hat{\StrainVector}\coloneq
\bm{\gamma}
+\bm{\kappa}\times\bm{r}
+\WarpTotalShapes\,\bm{\alpha}'
+\gradS (\WarpAxialShapes\cdot\bm{\alpha})
% \label{eq:linear-gl-strain-pi}
\end{SaveEquation}

\subparagraph{Section kinematic operator}
From \eqref{eq:linear-gl-strain-pi} the variation
\(
\delta\hat{\bm{\varepsilon}}=(\TotalStrainMatrix\,\delta\bm{e})
\)
furnishes:
\begin{equation}
\TotalStrainMatrix
=\left[
\;\mathbf{1}
~~~-\bm{r}^{\times}
~~~\FrameBasis\otimes\bm{\varphi}
~~~~\FrameBasis_{\beta}\otimes\nabla_{\beta}\WarpAxialShapes
\right],
\label{eq:def-section-kinematic-matrix-Phi}
\end{equation}
Thus, for any virtual increment,
\begin{equation}
\delta\StrainVector
=\delta\bm{\gamma}
-\bm{r}\times\delta\bm{\kappa}
+\WarpTotalShapes\,\delta\bm{\alpha}'
+\bm{\kappa}\times\WarpTotalShapes\,\delta\bm{\alpha}
+\gradS\big(\FrameBasis\cdot\bm{\Phi}\,\delta\bm{\alpha}\big),
\label{eq:delta-epsilon-Phi}
\end{equation}


\subsubsection{Resultants}

It is natural [TODO: elaborate on why; it follows equivalently from virtual work, stationary variational principles, and Galerkin weighted residuals] to impose the requirement:
\begin{equation}
\left<\TraceTotalEnergyStress,\,\delta\FrameStrain\right>_{\FrameBody}=\left<\StressTensor ,\,\delta\StrainTensor\right>_{\hat{\FrameBody}}
\label{eq:seSE}
\end{equation}
where \(\delta \bm{e}\) and \(\delta \StrainTensor\) are arbitrary variations of the strain variables.

Recall the Frobenius inner product 
\(
    \left.\bm{A} \cdot \bm{B}\right. \coloneq \operatorname{tr} \bm{A}\bm{B}^{\top}
\)
for two linear operators \(\bm{A}\) and \(\bm{B}\). 
%
If \(\delta \StrainTensor\) has the representation \(\left(\TotalStrainMatrix \delta\FrameStrain\right) \stackrel{s}{\otimes}\FrameBasis\) for some second-order tensor \(\TotalStrainMatrix\), then
\begin{equation*}
    \left<\StressTensor ,\,\delta\StrainTensor\right> = \SectionIntegral{ 
    \operatorname{tr} \left[\StressTensor\, \tfrac{1}{2}\left(\TotalStrainMatrix\delta \FrameStrain \otimes \FrameBasis + \FrameBasis\otimes \TotalStrainMatrix\delta \FrameStrain \right)
    \right]
    }
\end{equation*}
%
Exploiting linearity of the trace together with the identity \(\operatorname{tr} (\bm{a}\otimes\bm{b}) = \bm{a}\cdot \bm{b}\) for arbitrary vectors \(\bm{a}\) and \(\bm{b}\) yields
\begin{equation*}
    \left<\StressTensor ,\,\delta\StrainTensor\right> = \delta \FrameStrain\cdot \SectionIntegral{
        \TotalStrainMatrix^{\top}\tfrac{1}{2}\left(\StressTensor \,\FrameBasis + \StressTensor^{\top}\FrameBasis\right)
    }
\end{equation*}
and consequently by requiring \eqref{eq:seSE} the conjugate resultants \(\TraceTotalEnergyStress\) are obtained:
\begin{equation*}
\TraceTotalEnergyStress \equiv \SectionIntegral{\TotalStrainMatrix^{\top}\tfrac{1}{2}\left(\StressTensor\,\FrameBasis + \StressTensor^{\top}\FrameBasis\right)}
\end{equation*}
For a symmetric stress tensor \(\StressTensor\) one further has
\begin{equation}
\TraceTotalEnergyStress = \SectionIntegral{
  \TotalStrainMatrix^{\top} \StressVector
}
\label{eq:stress-resultants}
\end{equation}
where recall the stress vector \(\StressVector\) is defined by \eqref{eq:def-stress-vector}.

\begin{Details}
Note that for $\bm{\omega} \coloneq (\bm{\alpha}(\reflenx)\cdot\bm{\varphi}(y,z))\FrameBasis$
$$
\nabla \bm{\omega} = \nabla (\bm{\alpha}\cdot\bm{\varphi}\FrameBasis) = \FrameBasis\otimes\nabla(\bm{\alpha}\cdot\bm{\varphi}) = \FrameBasis\otimes\nabla\omega
$$

$$
\nabla \omega = (\nabla\bm{\alpha})^{\top}\bm{\varphi}  +(\nabla\bm{\varphi})^{\top}\bm{\alpha} 
$$
Using $\nabla \bm{\alpha} = \bm{\alpha}'\otimes\FrameBasis$
$$
\nabla \omega = (\bm{\alpha}^{\prime}\cdot\bm{\varphi}) \FrameBasis +(\nabla\bm{\varphi})^{\top}\bm{\alpha}
$$
and rearranging
$$
\nabla \omega = (\FrameBasis\otimes\bm{\varphi}) \bm{\alpha}' + (\nabla\bm{\varphi})^{\top}\bm{\alpha}
$$

$$
(\mathbf{a}\stackrel{\mathrm{s}}{\otimes}\FrameBasis)\FrameBasis
=\frac{1}{2} \left(\mathbf{a} + (\mathbf{a}\cdot\FrameBasis)\FrameBasis\right)
$$

$$
\delta \left((\mathbf{a}\stackrel{\mathrm{s}}{\otimes}\FrameBasis)\FrameBasis\right) = \frac{1}{2}\left(\delta \mathbf{a} + \delta a_1 \FrameBasis\right)
$$
\end{Details}


\subsection{Enhanced Strains}\label{sec:enhan-section}
The section strain field \(\StrainVector\) may be enhanced in a manner analogous to the three-field variational formulation of elasticity \citep{simo1986variational,simo1990class,simo1992geometrically}:
% \begin{equation*}
% \bm{E} = \hat{\bm{E}}(\hat{\bm{x}}) + \tilde{\bm{E}}
% \end{equation*}
\begin{SaveEquation}{eq:enhan-strain}
% \StrainVector(\reflenx,\bm{r}) = \Big(\RigidStrainMatrix  + \RigidStrainMatrix\RigidShift + \WarpStrainMatrix \ShapeShift\Big) \,\bm{e}(\reflenx)
\StrainVector(\reflenx,\bm{r}) =  \operatorname{A}[\bm{e}] + \EnhanShapes(\bm{r}) \bm{\EnhanParam}(x)
\end{SaveEquation}
where \(\hat{\bm{E}} = \hat{\bm{\varepsilon}}\symbun\FrameBasis\) is the compatible strain tensor derived from the displacement field \eqref{eq:finite-alignment-field}, and $\tilde{\bm{E}}$ is the enhanced part of the strain field. 
% This may be decomposed as
% \begin{equation*}
% \tilde{\bm{E}} = \tilde{\bm{\varepsilon}} \symbun \FrameBasis + \StrainTensorJ
% \end{equation*}
% where \(\StrainTensorJ\) is the part of the strain tensor with components on \(\PlaneBasis_{\alpha}\otimes \PlaneBasis_{\beta}\), which is solved for by enforcing the stress constraints:
% \[
% \PlaneBasis_{\alpha} \cdot \StressTensor \PlaneBasis_{\beta} = 0
% \]
% Our interest is in the enhanced strain vector \(\tilde{\StrainVector}\). 
Let $\operatorname{B}^*\left[\mathcal{E}\right] \subset \mathscr{E}$ be the space of compatible linear strain tensors generated by the space $\mathcal{E}$ of rod strain fields \(\bm{e}\); i.e. 
$$
\begin{aligned}
\operatorname{B^*}\left[\mathcal{E}\right]
% &=\left\{\bm{E}^h=\sum_{e=1}^{n_{\mathrm{clm}}} \chi_e \bm{\varepsilon}_e(\xi) \symbun\FrameBasis
% \; \bigm| 
% \bm{\varepsilon}_e(\xi)=\sum_{I=1}^{n_{\text {node }}^e} \mathbf{u}_I^e \otimes \operatorname{B^*}[\bm{e}^h(\xi)]
% \right\} \\
&\coloneq \left\{
\hat{\StrainVector} \symbun \FrameBasis \in \mathscr{E} 
\bigm| 
%\hat{\StrainVector}^h=\sum_{e=1}^{n_{\text {elm }}} \hat{\StrainVector}_e(\xi) \chi_e, \text { with } 
\hat{\StrainVector}%_e(\xi)
=\TotalStrainMatrix(\xi,\hat{\rho}) \bm{e},
\quad \bm{e} \in \mathcal{E}
\right\}
\end{aligned}
$$


% in the context of nonlinear solid mechanics, a body B is embedded in an ambient space A by a map $\chi \in C$, where $C$ is an infinite dimensional manifold of maps $\{\chi: B \rightarrow A\}$. 
% In the abscence of constraints (eg, shell and rod theories) $C$ may generally be regarded as a (Linear) Hilbert or Sobolev space. 
% The space of spatial strain measures may be constructed: $\{E: \chi(B)->\chi(B)...\}$.

% \begin{equation}
% \Pi(\mathbf{u}, \tilde{\bm{E}}, \bm{T}) 
% :=\int_0\left.W\left(\bm{x}, \bm{E}_x+\tilde{\bm{\varepsilon}}\right)-\bm{T} \cdot \tilde{\bm{E}}
% \right. \,\mathrm{d} V
% -\int_0 \rho_0 \bm{b} \cdot \mathbf{u} \mathrm{~d} V-\int \tilde{\mathbf{t}} \cdot \mathbf{u} \, \mathrm{d} \Gamma
% \end{equation}

Define the finite dimensional space generated by assumed strain interpolations \(\EnhanShapes\bigl(\xi,\hat{\rho}\bigr)\)
%subspace $\tilde{\mathscr{E}}^h \subset \tilde{\mathscr{E}}$ by the expression
% \begin{equation}
% \tilde{\mathscr{E}}^h
% =\left\{\tilde{\bm{H}}^h \in \tilde{H}
% \;: \tilde{\bm{H}}^h=\sum_{e=1}^{n_{e \mathrm{~lm}}} \tilde{\bm{H}}_e(\xi) \chi_e ; 
% \;
% \tilde{\bm{H}}_e(\xi)=\sum_{I=1} \tilde{\alpha}_I \otimes \mathbf{G}_e^I \text { for } \tilde{\alpha}_I \in \mathbb{R}^{n_{\mathrm{dim}}}
% \right\}
% \label{eq:enhanced-space}
% \end{equation}
\begin{equation*}
\tilde{\mathscr{E}}
\coloneq 
\left\{
\tilde{\StrainVector} \symbun \FrameBasis \in \tilde{\mathscr{E}} \mid %\tilde{\StrainVector}^h=\sum_{e=1}^{n_{\text {lim }}} \tilde{\gamma}_e(\xi) \chi_e, 
%\text { with } 
\tilde{\bm{\varepsilon}}(\xi,\hat{\rho})=\EnhanShapes(\xi, \hat{\rho}) \bm{\EnhanParam} , \text { and } \bm{\EnhanParam} \in \mathbb{R}^{\EnhanDim}\right\}
\label{eq:SR16}
\end{equation*}
where $\tilde{\bm{\alpha}} \in \mathbb{R}^{n_a}$ are $\EnhanDim$ local parameters. %, 
% and $\mathbf{G}^I: \PlaneManifold \rightarrow \mathbb{R}^{n_{\mathrm{dim}}}$ are $\EnhanDim$ vectors of prescribed interpolation functions. 
%
% Here, ${\tilde{\bm{\alpha}}}_{\xi}$ are $\EnhanDim$ section strain parameters, and $\EnhanShapes (\PlanePosition)$ is a matrix of $\mathbb{R}^{n_{n+r}} \times \mathbb{R}^{n_\tau}$ prescribed functions with linearly independent columns, which define the enhanced strain interpolation. 
Consider conditions analogous to those of \cite{simo1992geometrically}:
\begin{enumerate}[label=\roman*]
    \item Stability. 
    % The subspaces spanned by the strains $\operatorname{B^*} [N^I]$ and the assumed functions $\mathbf{G}_e^I$ in \eqref{eq:SR16} have null intersection. 
%
    The enhanced strain space $\tilde{\mathscr{E}}$ and the compatible strain space \(\operatorname{B^*}\left[\mathcal{E}\right]\) have null intersection
    
    $$
    \tilde{\mathscr{E}} \cap \operatorname{B^*}\left[\mathcal{E}\right]=\varnothing
    $$


\item Consistency. % $S^h$ and $\tilde{\mathscr{E}}^h$ are $L_2$-orthogonal. 
    The subspace of admissible stresses, denoted by $\StressSpaceH \subset L$, is $L_2$ -orthogonal to the space $\tilde{\mathscr{E}}$ of enhanced strain fields. 
    This requires \(\residn(\bm{\EnhanParam}) = \mathbf{0}\) with
    \begin{SaveEquation}{eq:def-enhan-residual}
    \residn(\bm{\EnhanParam}) \coloneq \int_{\Section} \EnhanShapes^{\top} \StressVector \left(\operatorname{A}[\bm{e}] + \EnhanShapes \bm{\EnhanParam} \right) \dA
    % \label{eq:def-enhan-residual}
    \end{SaveEquation}
% \item Patch test. The space $\StressSpaceH$ of admissible nominal stress fields contains piece-wise constant functions. 
% In connection with \eqref{eq:def-enhan-residual}, this furnishes the condition:
\end{enumerate}
Condition (i) is essential. Condition (ii) is ideal and justified as follows:
\begin{itemize}
    \item When...
\end{itemize}
% In traditional assumed strain formulations a third criteria is required
% which asserts that the space of admissible nominal stress fields contains piece-wise constant functions, and implies \citep{simo1990class}:
% \begin{equation}
% \int_{\Section \times \FrameBody} \EnhanShapes \dA \, \dd \xi = \mathbf{0}
% \label{eq:condiii}
% \end{equation}
% However, such a condition would not make sense in the present context. 


\subsection{Linearization}

Linearizing \eqref{eq:stress-resultants} furnishes:
\begin{equation}
    \mathcal{L}\bm{s} = \Ke \dd \bm{e} + \Ken \dd \bm{\eta}
\label{eq:Lins}
\end{equation}
where:
\begin{subequations}
\begin{align}
\Ke &\coloneq \int \TotalStrainMatrix^{\top}\mathbb{C}\TotalStrainMatrix \dA  + \Kwg \label{eq:Kee} \\
\Ken &\coloneq \int \TotalStrainMatrix^{\top}\mathbb{C}\EnhanShapes\dA \label{eq:Ken} \\
\end{align}
\end{subequations}
and \(\Kwg\) is obtained by linearizing the Wagner term \eqref{eq:wagner-gl-strain-pi}:
\begin{equation}
\begin{aligned}
\Kwg =
\Twist & \SectionIntegral{
\left[\begin{matrix}
\mathbf{0} & r^{2}  \mathbb{C} \FrameBasis \otimes \FrameBasis & \mathbf{0} & \mathbf{0} \\
r^{2}  \FrameBasis \otimes \FrameBasis \mathbb{C} 
& r^{2}  \left(\bm{r}^{\times} \mathbb{C} \FrameBasis \otimes \FrameBasis - \FrameBasis \otimes \FrameBasis \mathbb{C} \bm{r}^{\times}\right) & r^{2}  \FrameBasis \otimes \FrameBasis \mathbb{C} \FrameBasis\otimes \bm{\varphi} & r^{2}  \FrameBasis \otimes \FrameBasis \mathbb{C} \FrameBasis_{\beta} \otimes \nabla_{\beta} \bm{\varphi}\\
\mathbf{0} & r^{2}  \bm{\varphi} \otimes \FrameBasis \mathbb{C} \FrameBasis \otimes \FrameBasis & \mathbf{0} & \mathbf{0} \\
\mathbf{0} & r^{2}  \nabla_{\beta} \bm{\varphi} \otimes \FrameBasis_{\beta} \mathbb{C} \FrameBasis \otimes \FrameBasis & \mathbf{0} & \mathbf{0} 
\end{matrix}\right]} \\
+\Twist^2 & \SectionIntegral{\left[\begin{matrix}
\mathbf{0} & \mathbf{0} & \mathbf{0} & \mathbf{0} \\
\mathbf{0} & r^{4} \FrameBasis \otimes \FrameBasis \mathbb{C} \FrameBasis \otimes \FrameBasis & \mathbf{0} & \mathbf{0} \\
\mathbf{0} & \mathbf{0} & \mathbf{0} & \mathbf{0} \\
\mathbf{0} & \mathbf{0} & \mathbf{0} & \mathbf{0}
\end{matrix}\right]
} \\
&+\SectionIntegral{\left[\begin{matrix}
\mathbf{0} & \mathbf{0} & \mathbf{0} & \mathbf{0} \\
\mathbf{0} & r^2 \AxialStress & \mathbf{0} & \mathbf{0} \\
\mathbf{0} & \mathbf{0} & \mathbf{0} & \mathbf{0} \\
\mathbf{0} & \mathbf{0} & \mathbf{0} & \mathbf{0} 
\end{matrix}\right]
}
\end{aligned}
\label{eq:section-wagner-tangent}
\end{equation}

Linearizing \eqref{eq:def-enhan-residual} furnishes:
\begin{equation}
\mathcal{L}\residn = \mathbf{K}_{e\eta}^{\top}\dd \bm{e} + \mathbf{K}_{\eta \eta}\dd \bm{\eta}
\qquad\implies\qquad 
\dd\bm{\eta} = \widetilde{\mathbf{K}} \dd \bm{e}
\label{eq:Linn}
\end{equation}
with
\begin{subequations}
\begin{align}
% \mathbf{K}_{\eta e} &\coloneq \int \EnhanShapes^{\top}\mathbb{C}\TotalStrainMatrix\dA \\
    \Knn &\coloneq \int \EnhanShapes^{\top}\mathbb{C}\EnhanShapes\dA \label{eq:Knn} \\
    \Kn &\coloneq -\Knn^{-1}\Ken^{\top} \label{eq:Kn}
\end{align}
\end{subequations}

From \eqref{eq:Lins} and \eqref{eq:Linn}:
\begin{SaveEquation}{eq:section-tangent}
\Ks = \Ke - \mathbf{K}_{\eta e}^{\top} \Kn
\end{SaveEquation}
